% 6. International climate policy: majority support even when people aware of the costs
%===============================================================================================
% - burden-sharing: every country should participate, burden should be equitably shared
%===============================================================================================

International redistributive mechanisms, including some climate-policy schemes, must determine how countries share the costs. Public acceptance is likely to be higher when a broad coalition participates and contributions reflect each country's capacity.

Survey experiments on climate burden-sharing confirm the importance of institutional design. Support increases when many countries participate and share the costs rather than leaving most of the burden to the respondent's country. Respondents nevertheless remain willing to accept substantial domestic contributions. \citet{bechtel_mass_2013} conducted survey experiments in 2012 with nationally representative samples from France, Germany, the United Kingdom, and the United States. Each respondent completed four comparisons between two randomly generated international climate agreements. For each comparison, respondents selected their preferred agreement and reported how likely they were to vote for or against their country's participation in a referendum. The authors estimate that moving from an agreement resembling the one then under discussion to the design predicted to maximize support would raise the share voting in favor from 42\% to 60\% in France, from 37\% to 60\% in Germany, from 36\% to 51\% in the United Kingdom, and from 29\% to 47\% in the United States \citep[Table S3]{bechtel_mass_2013}.


The agreements varied in their monthly household cost, the distribution of costs across countries, the number of participating countries, the share of global emissions covered, the monitoring institution, and the sanctions for non-compliance. These design features strongly affected support. Increasing costs from 0.5\% to 1\% of national GDP reduced the probability of selecting an agreement by approximately 10 percentage points. By contrast, increasing participation from 20 to 160 countries out of 192 raised it by about 15 points \citep[Table 1 \& Figure 2]{bechtel_mass_2013}.

Focusing on North--South climate finance, \citet{gampfer_obtaining_2014} conducted survey experiments with nationally representative samples in the United States and Germany between October and November 2013. Respondents considered six pairs of hypothetical climate-finance mechanisms and indicated, for each pair, ``which of the two proposals they preferred''. The mechanisms varied in recipient countries' climate damages, income, emissions, and government efficiency; the actors controlling allocation decisions; the use of funds for mitigation, adaptation, or both; and the distribution of contributions between the respondent's country and other donor countries. Donor burden-sharing produced by far the largest effect. Relative to a mechanism in which other countries provided 10\% of the funds and the respondent's country provided 90\%, reversing these shares increased the probability of selection by 28 percentage points in the United States and 23 points in Germany. Respondents nevertheless accepted substantial domestic contributions: most of the increase in support occurred when the share paid by other countries rose from 10\% to 30\% or 50\%. 
When decisions about which projects would receive funding were made jointly by contributing and recipient countries rather than by recipient countries alone, the probability of selecting the mechanism increased by 6 percentage points in the United States and 9 points in Germany. Respondents also preferred mechanisms that financed both mitigation and adaptation to those that financed adaptation alone \citep[Table 1 \& Figure 2]{gampfer_obtaining_2014}.

%===============================================================================================
% - Global Climate Scheme
%===============================================================================================

\citet{fabre_public_2026} examines public acceptance of a Global Climate Scheme (GCS). The description shown to respondents began: ``In 2015, all countries agreed to contain global warming `well below +2\textdegree{}C'.'' The proposed scheme would impose a global ceiling on greenhouse-gas emissions and require polluting firms to purchase permits within that limit. By incorporating the cost of emissions into fossil-fuel prices, the mechanism was expected to reduce fossil-fuel consumption and, consequently, greenhouse-gas emissions. Permit revenues would finance equal monthly payments to adults worldwide, reflecting the stated principle that ``each human has an equal right to pollute.'' Respondents were informed that an average transfer of $35 per month would lift approximately 600 million people earning less than \$2 per day out of extreme poverty. They were also shown country-specific estimates of the net monthly loss borne by a typical resident and of the corresponding transfer: \$90 and \$35 in the United States, \texteuro 15 and \texteuro 20 in France, \texteuro 45 and \texteuro 20 in Germany, CHF 15 and CHF 20 in Switzerland, and RUB 2,500 and RUB 3,000 in Russia. In the United States, the projected 2\% increase in consumer prices implied a net monthly loss of \$90 after subtracting the $35 transfer \citep[Table S2]{fabre_public_2026}.

%LB: il faudra juger de l'utilité de ce paragraphe --> la Figure S30 suit ce paragraphe.
The GCS is accepted by 55\% of respondents in the pooled sample. Acceptance is lowest in the United States and Russia, at 49\%. It reaches 61\% in Europe as a whole, including 65\% in France, 53\% in Germany, and 58\% in the United Kingdom. Among the eleven surveyed countries, Saudi Arabia displays the highest acceptance rate, at 85\% \citep[Figures 6 \& S30]{fabre_public_2026}.


%===============================================================================================
% - International Climate Schemes
%===============================================================================================

The International Climate Scheme (ICS) applies the same cap-and-trade mechanism as the GCS, but within different coalitions of countries. Respondents are randomly assigned to one of four variants and shown a map displaying ``a possible set of countries'' that would participate in the scheme, together with their share of global emissions. In the \textit{Low} variant, the coalition comprises the Global South and the European Union, but excludes China, covering 33\% of global emissions. The \textit{Mid} variant replaces the EU with China and brings together China and the Global South, which account for 56\% of global emissions. The \textit{High} variant adds the EU, the United Kingdom, Japan, Canada, South Korea, and other high-income countries to the \textit{Mid} coalition, raising coverage to 72\% of global emissions. Finally, the \textit{High color} variant retains the same coalition as the \textit{High} variant but uses a colored map to display the net financial gain or cost of the scheme for each country. Respondents assigned to this variant were told that ``a provision would prevent the Global Climate Scheme from harming low- and middle-income countries'', so countries such as China ``would neither win nor lose financially'' \citep[Questions 31--35 and Figure 5]{fabre_public_2026}. Wider participation increases acceptance only modestly: support rises from 65\% in the \textit{Low} variant to 68\% in the \textit{High} variant. By contrast, explicitly displaying national gains and losses reduces support by 6.6 percentage points relative to the \textit{High} variant, suggesting that the salience of distributive consequences matters more than the precise size of the coalition \citep[Figure 6]{fabre_public_2026}.

All four ICS variants attract majority support in every country surveyed \citep[Figures 6 \& S30]{fabre_public_2026}. Acceptance of the \textit{Low} coverage variant is similar to that of the National Climate Scheme, although the ICS is more costly for the average respondent. As Fabre notes, this comparison ``suggests that the average respondent is willing to pay the ICS's higher cost'' in exchange for guaranteed poverty alleviation and decarbonization in the Global South. The ICS also receives greater acceptance than the GCS, but this may reflect the greater credibility and precision of a proposal that lists participating countries, its visual presentation, or the lower salience of individual costs. Nevertheless, ``the data does not allow testing these different hypotheses'' \citep{fabre_public_2026}.

\begin{figure}[htbp]
    \centering
    \includegraphics[
        page=76,
        width=\textwidth,
        trim=25 90 20 470,
        clip
    ]{figures/fabre_2026.pdf}
    \caption{Support for the National, Global, and International Climate Schemes by country \citep[Figure S30]{fabre_public_2026}.}
    \label{fig:ics_support_country}
\end{figure}


%===============================================================================================
% - climate policy supported if it's perceived as fair, environmentally effective and in one's interest, suggesting majority support as soon as two out of three elements are there (Douenne & Fabre 2022); if int'l climate policy seen as fair and effective, logical that it's supported even if against one's interest
%===============================================================================================

%LB: paragraphe raccourci, j'ai quand même gardé deux chiffres.
\citet{douenne_yellow_2022} show that acceptance of climate policy depends strongly on three perceived features: whether the policy is fair, environmentally effective, and in the respondent's own financial interest. In their estimates, believing that one would not lose financially increases acceptance by 53 percentage points, while believing that the policy is environmentally effective raises approval by 42 points. The authors therefore suggest that convincing citizens on more than one of these dimensions would likely generate large majority support, although they cannot causally identify their combined effect. Accordingly, an international climate policy may remain majoritarian even when it imposes a personal cost, provided that it is regarded as both fair and environmentally effective \citep[Tables 6 \& 8]{douenne_yellow_2022}.
